\documentclass[12pt,a4paper]{article}

%% PACHETE MATEMATICE
\usepackage{amsmath,amssymb,amsfonts}
\usepackage{amsthm}
\usepackage{mathpazo}
\usepackage{eulervm}
\usepackage{enumerate}
\usepackage{color}
\usepackage{graphicx}
\usepackage[all]{xy}
\usepackage{xcolor}
\usepackage{framed}
\usepackage[unicode]{hyperref}
\setcounter{MaxMatrixCols}{20}
\usepackage{multicol}
%\usepackage[romanian]{babel}


%% ASPECT
\linespread{1.05}
\usepackage[T1]{fontenc}
\usepackage[utf8x]{inputenc}
\usepackage[margin=2cm]{geometry}
% \usepackage[color]{showkeys}
\usepackage{anyfontsize}
\usepackage{titlesec}
\titleformat*{\section}{\large\bfseries}

\usepackage{fancyhdr}
\pagestyle{fancy}
\rhead{\small \color{gray} Notițe de Adrian Manea}
\lhead{\small \color{gray} IntProb-FIA, 2017---2018, L. Lemnete \& M. Olteanu}


\usepackage[autostyle = false, style = german]{csquotes}
\usepackage[romanian]{babel}
\newcommand\qq{\enquote}

%% COMENZILE MELE
\renewcommand*{\proofname}{Dem.:}
\newcommand\bb{\mathbb}
\newcommand\dr{\mathrm}
\newcommand\kal{\mathcal}
\newcommand\fk{\mathfrak}
\newcommand\op{\oplus}
\newcommand\ot{\otimes}
\newcommand\ds{\displaystyle}
\newcommand\id{\indent}
\newcommand\nid{\noindent}
\newcommand\seq{\subseteq}
\newcommand\sneq{\subsetneq}
\newcommand\speq{\supseteq}
\newcommand\spneq{\supsetneq}
\newcommand\rar{\rightarrow}
\newcommand\Rar{\Rightarrow}
\newcommand\xrar{\xrightarrow}
\newcommand\lar{\leftarrow}
\newcommand\Lar{\Leftarrow}
\newcommand\xlar{\xleftarrow}
\newcommand\lrar{\leftrightarrow}
\newcommand\Lrar{\Leftrightarrow}
\newcommand\ideal{\trianglelefteq}
\newcommand\ai{\text{ a.\^{i}. }}
\newcommand\sk{(\textit{Schi\c{t}\u{a}}) }
\newcommand\rhup{\rightharpoonup}
\newcommand\rhdn{\rightharpoondown}
\newcommand\lhup{\leftharpoonup}
\newcommand\lhdn{\leftharpoondown}
\newcommand\vid{\emptyset}
\newcommand\wh{\widehat}
\newcommand\ol{\overline}
\newcommand\NN{\mathbb{N}}
\newcommand\RR{\mathbb{R}}
\newcommand\ZZ{\mathbb{Z}}
\newcommand\QQ{\mathbb{Q}}
\newcommand\CC{\mathbb{C}}

%% STILURILE MELE
\newtheoremstyle{dotless}{}{}{\itshape}{}{\bfseries}{:}{ }{}

\theoremstyle{dotless}
\newtheorem{theorem}{Teorem\u{a}}[section]
\newtheorem{proposition}{Propozi\c{t}ie}[section]
\newtheorem{lemma}{Lem\u{a}}[section]
\newtheorem{corollary}{Corolar}[section]
\newtheorem{exercise}{Exerci\c{t}iu}

\newtheoremstyle{dotlessdr}{}{}{}{}{\bfseries}{:}{ }{}
\theoremstyle{dotlessdr}
\newtheorem{example}{Exemplu}[section]
\newtheorem{definition}{Defini\c{t}ie}[section]
\newtheorem{remark}{Observa\c{t}ie}[section]




\begin{document}

\begin{center}
  {\large\textbf{Seminar 10 \\ Formule integrale}}
\end{center}

\vspace{2cm}


\section{Formula Green-Riemann}
\label{sec:g-r}

Fie $(K, \partial K)$ un compact cu bord orientat inclus în $\RR^2$ și considerăm
o 1-formă diferențială de clasă $\kal{C}^1$ pe o vecinătate a lui $K$. Atunci
are loc \emph{formula Green-Riemann}:
\[
  \int_{\partial K} Pdx + Qdy = \iint_K \Big(\frac{\partial Q}{\partial x} - %
  \frac{\partial P}{\partial y} \Big) dx dy.
\]

O consecință imediată este o formulă pentru arie:
\[
  \mathcal{A}(K) = \frac{1}{2} \int_{\partial K} xdy - ydx.
\]

%%%%%%%%%%%%%%%%%%%%%%%%%%%%%%%%%%%%%%%%%%%%%%%%%%%%%%%%%%%%%%%%%%%%%%

\section{Formula Gauss-Ostrogradski}
\label{sec:g-o}

Considerăm $K$ o mulțime compactă, cu bord orientat după normala exterioară. Atunci,
pentru orice 2-formă diferențială:
\[
  \omega = Pdy \wedge dz + Q dz \wedge dx + R dx \wedge dy
\]
de clasă $\kal{C}^1$ pe o vecinătate a lui $K$ are loc egalitatea:
\[
  \int_{\partial K} P dy \wedge dz + Q dz \wedge dx + R dx \wedge dy = %
  \iiint_K \frac{\partial P}{\partial x} + \frac{\partial Q}{\partial y} + %
  \frac{\partial R}{\partial z} dxdydz.
\]

În notație vectorială, dacă $\vec{V} = P \vec{i} + Q \vec{j} + R \vec{k}$ este
cîmpul vectorial asociat 2-formei diferențiale $\omega$, atunci formula de mai sus
poate fi scrisă:
\[
  \int_{\partial K} \vec{V} \cdot \vec{n} d \sigma = \iiint_K \mathrm{div}\vec{V} dxdydz,
\]
unde $\vec{n}$ este normala exterioară la $\partial K$. În membrul stîng avem
fluxul cîmpului $\vec{V}$ prin suprafața $\partial K$, motiv pentru care formula
Gauss-Ostrogradski se mai numește \emph{formula flux-divergență}.

%%%%%%%%%%%%%%%%%%%%%%%%%%%%%%%%%%%%%%%%%%%%%%%%%%%%%%%%%%%%%%%%%%%%%%

\section{Formula lui Stokes}
\label{sec:stokes}

Fie $\Sigma$ o suprafață cu bord, orientată și fie:
\[
  \alpha = Pdx + Qdy + Rdz
\]
o 1-formă diferențială de clasă $\kal{C}^1$ pe o vecinătate a lui $\Sigma$.
Atunci are loc formula:
\begin{align*}
  \int_{\partial \Sigma} Pdx + Qdy + Rdz &= \int_\sigma \Big( \frac{\partial R}{\partial y} -
                                           \frac{\partial Q}{\partial z} \Big) dy \wedge dz +
                                           \Big( \frac{\partial P}{\partial z} -
                                           \frac{\partial R}{\partial x} \Big) dz \wedge dx +
                                           \Big( \frac{\partial Q}{\partial x} -
                                           \frac{\partial P}{\partial y}\Big) dx \wedge dy.
\end{align*}

În notație vectorială, dacă $\vec{V} = P \vec{i} + Q \vec{j} + R \vec{k}$ este cîmpul
vectorial asociat formei diferențiale $\alpha$, atunci formula lui Stokes se scrie:
\[
  \int_{\partial \Sigma} \vec{V} \cdot d\vec{r} = \int_\Sigma (\nabla \times \vec{V}) %
  \cdot \vec{n} d \sigma,
\]
unde $\nabla \times \vec{V} = \mathrm{rot}\vec{V}$ se numește \emph{rotorul} cîmpului
vectorial $\vec{V}$, calculat cu ajutorul produsului vectorial.

%%%%%%%%%%%%%%%%%%%%%%%%%%%%%%%%%%%%%%%%%%%%%%%%%%%%%%%%%%%%%%%%%%%%%%

\section{Exerciții}
\label{sec:exc-f}

1. Să se calculeze direct și folosind formula Green-Riemann integrala curbilinie
$\int_\Gamma \alpha$ în următoarele cazuri:
\begin{enumerate}[(a)]
\item $\alpha = y^2 dx + xdy$, unde $\Gamma$ este pătratul cu vîrfurile
  $A(0, 0), B(2, 0), C(2,2), D(0, 2)$;
\item $\alpha = ydx + x^2 dy$, unde $\Gamma$ este cercul cu centrul în origine și
  rază 2;
\item $\alpha = ydx - xdy$, unde $\Gamma$ este elipsa de semiaxe $a$ și $b$ și de
  centru $O$.
\end{enumerate}

\vspace{1cm}

2. Să se calculeze integralele, direct, apoi aplicînd formula Green-Riemann:
\begin{enumerate}[(a)]
\item $\ds\int_\Gamma e^{\frac{x^2}{a^2} + \frac{y^2}{b^2}} (-ydx + xdy)$, unde
  $\Gamma = \{(x, y) \mid \frac{x^2}{a^2} + \frac{y^2}{b^2} = 1 \}$;
\item $\ds\int_\Gamma xy dx + \frac{x^2}{2} dy$, unde $\Gamma$ este obținută prin:
  \[
    \Gamma = \{x^2 + y^2 = 1, x \leq 0 \leq y\} \cup \{x + y = -1, x,y \leq 0 \}.
  \]
\end{enumerate}

\vspace{1cm}

3. Să se calculeze circulația cîmpului vectorial $\vec{V}$ pe curba $\Gamma$ în
cazurile:
\begin{enumerate}[(a)]
\item $\vec{V} = y^2 \vec{i} + xy\vec{j}$, unde:
  \[
    \Gamma = \{x^2 + y^2 = 1, y > 0 \} \cup \{ y = x^2 - 1, y \leq 0 \};
  \]
\item $\vec{V} = e^x \cos y \vec{i} - e^x \sin y \vec{j}$, unde $\Gamma$ este
  o curbă arbitrară din semiplanul superior, ce unește punctele $A(1, 0), B(-1, 0)$,
  cu sensul de la $A$ către $B$.
\end{enumerate}

\vspace{1cm}

4. Să se calculeze integrala de suprafață $\int_\Sigma \omega$ în următoarele cazuri:
\begin{enumerate}[(a)]
\item $\omega = x^2 dy \wedge dz - 2xy dz \wedge dx + z^3 dx \wedge dy$, iar mulțimea
  $\Sigma = \{x^2 + y^2 + z^2 = 9 \}$;
\item $\omega = yz dy \wedge dz - (x + z) dz \wedge dx + (x^2 + y^2 + 3z) dx \wedge dy$,
  iar mulțimea:
  \[
    \Sigma = \{ x^2 + y^2 = 4 - 2z, z \geq 1 \} \cup \{x^2 + y^2 \leq 4 - 2z, z = 1 \};
  \]
\item $\omega = x(z + 3) dy \wedge dz + yz dz \wedge dx - (z + z^2) dx \wedge dy$,
  iar mulțimea:
  \[
    \Sigma = \{ x^2 + y^2 + z^2 = 1, z \geq 0\}.
  \]
\end{enumerate}

\vspace{1cm}

5. Să se calculeze, folosind formula lui Stokes, integrala curbilinie
$\int_\Gamma \alpha$, în următoarele cazuri:
\begin{enumerate}[(a)]
\item $\alpha = (y - z) dx + (z - x)dy + (x - y) dz$, iar mulțimea $\Gamma: z = x^2 + y^2, z = 1$;
\item $\alpha = ydx + zdy + xdz$, iar mulțimea:
  \[
    \Gamma: x^2 + y^2 + z^2 = 1, \quad x + y + z = 0.
  \]
\end{enumerate}

\vspace{1cm}

6. Să se calculeze circulația cîmpului vectorial:
\[
  \vec{V} = (y^2 + z^2)\vec{i} + (x^2 + z^2)\vec{j} + (x^2 + y^2)\vec{k},
\]
pe curba $\Gamma: x^2 + y^2 + z^2 = R^2, ax + by + cz = 0$.

\end{document}